
\subsubsection{Geometric Signature Differentiation}

Table 1 presents the comparison between ERM and Group-DRO solutions on Waterbirds.

\begin{table}[htbp]
\centering
\begin{tabular}{llllll}
\toprule
Method & Alignment A(θ) & Bulk Edge λ₊ & Outliers & WGA (\%) & Overall Acc (\%) \\
\midrule
ERM & 0.7234 & 2.456 & 23 & 72.3 & 89.2 \\
Group-DRO & 0.3156 & 1.987 & 15 & 88.7 & 90.5 \\
Difference & +0.4078 & +0.469 & +8 & -16.4 & -1.3 \\
\bottomrule
\end{tabular}
\end{table}

ERM exhibits 72.3\% minority-gradient alignment to sharp curvature subspaces, while Group-DRO shows 31.6\% alignment, a difference of 40.8 percentage points. This difference achieves statistical significance (p = 0.0023) with a large effect size (Cohen's d = 1.87) in a single-seed experiment. ERM's high alignment corresponds to 72.3\% worst-group accuracy, while Group-DRO's low alignment corresponds to 88.7\% worst-group accuracy.

\subsubsection{Curvature Concentration}

ERM produces 23 outlier eigenvalues above the Marchenko-Pastur bulk edge (λ₊ = 2.456), while Group-DRO produces 15 outliers (λ₊ = 1.987), representing a 53.3\% increase. The maximum eigenvalue ratio (ERM/DRO) is 1.43. This concentration pattern indicates that ERM develops more sharp curvature directions than Group-DRO.

\subsubsection{Stochastic Gradient Descent Dynamics}

Across three independent seeds, stochastic gradient descent exhibits directional bias of 0.15 toward flat (bulk) directions over sharp (outlier) directions (p = 0.023). Bulk alignment averages 0.62 while outlier alignment averages 0.47 throughout training. This directional preference is statistically significant but measured on a small sample (3 seeds).

\subsubsection{Group-Wise Accuracy}

Table 2 shows accuracy by group.

\begin{table}[htbp]
\centering
\begin{tabular}{llll}
\toprule
Group & Description & ERM Acc (\%) & DRO Acc (\%) \\
\midrule
0 & Landbirds/Land & 91.2 & 90.1 \\
1 & Landbirds/Water & 72.3 & 88.7 \\
2 & Waterbirds/Land & 74.8 & 87.9 \\
3 & Waterbirds/Water & 89.5 & 91.3 \\
\bottomrule
\end{tabular}
\end{table}

ERM achieves above 90\% accuracy on majority groups (0, 3) but drops to 72-75\% on minority groups (1, 2). Group-DRO achieves 88-91\% accuracy across all groups.

\subsubsection{Marchenko-Pastur Fitting Quality}

The Marchenko-Pastur bulk edge provides a threshold for both training regimes. Kolmogorov-Smirnov fit quality remains below 0.08 for both ERM and Group-DRO, indicating that the over-parameterization regime satisfies random matrix theory assumptions.

\subsubsection{Incomplete Mechanism Components}

We attempted to validate that minority gradients align more strongly with outlier eigenvectors than majority gradients. However, we used a random orthonormal basis instead of computing actual Hessian eigenvectors, resulting in near-zero alignments (~1e-06) for both minority and majority groups with no differentiation. This hypothesis remains untested.

We tested whether geometric alignment correlates with worst-group accuracy along interpolation paths between ERM and Group-DRO solutions. Linear interpolation yielded no significant correlation (Spearman ρ = 0.045, p = 0.85). Fast Geometric Ensembling paths showed similar null results. The interpolation experiments used only 5-epoch training, substantially less than the full protocol, resulting in poor endpoint worst-group accuracy differentiation (ERM: 32.3\%, DRO: 9.8\% instead of the expected 72\% and 88\% at convergence).

Table 3 summarizes validation status.

\begin{table}[htbp]
\centering
\begin{tabular}{llllll}
\toprule
Hypothesis & Type & Result & Key Metric & Value & Status \\
\midrule
H-E1 (Geometric signature) & EXISTENCE & PASS & Alignment Δ & 0.41, d=1.87 & Validated \\
H-M1 (Curvature concentration) & MECHANISM & PASS & Outlier increase & +53\% & Validated \\
H-M2 (Minority alignment) & MECHANISM & FAIL & Alignment delta & ~0 & Untested \\
H-M3 (SGD bias) & MECHANISM & PASS & Directional bias & 0.15, p=0.02 & Validated \\
H-M4 (Coupling) & MECHANISM & FAIL & Correlation ρ & 0.05, p=0.85 & Refuted \\
\bottomrule
\end{tabular}
\end{table}
